\section{Verified DFS interlude}
\label{sec:dfs}

The basic pebble game of \cref{sec:pebble-game} repeatedly searches a
partial orientation for a vertex with a free pebble reachable from a
target endpoint, returning the witness path so it can be reversed.
The inner primitive is a depth-first search on a directed graph with
predicate-based early termination; we formalize it once here, as a
short interlude, and consume it as a black box in
\cref{def:tryReachPebble}.

Mathlib's \texttt{Combinatorics.SimpleGraph}
$\mathrm{Walk}$ / $\mathrm{Reachable}$ infrastructure is declarative
(``a walk exists'', ``a reverse exists''), not algorithmic. We model
termination on \texttt{Batteries.UnionFind}, whose
\texttt{termination\_by self.rankMax - self.rank x} is a
strictly-decreasing numeric measure on a finite data structure;
here the measure is specialised to a $\mathrm{Finset}$-cardinality
on unvisited vertices.

\begin{definition}[Directed walk]
  \label{def:directed-walk}
  \lean{CombinatorialRigidity.Search.DirectedWalk,
    CombinatorialRigidity.Search.DirectedWalk.IsPath}
  \leanok
  Let $V$ be a finite type with decidable equality and $R : V \to V
  \to \mathrm{Prop}$ a relation. A \emph{directed $R$-walk} from
  $u$ to $w$ is a non-empty list $u = v_0, v_1, \ldots, v_m = w$ of
  vertices satisfying $R\,v_i\,v_{i+1}$ for $0 \le i < m$. The walk
  is \emph{simple} if all $v_i$ are distinct.
\end{definition}

\begin{definition}[Reachable-finding DFS]
  \label{def:reachable-finding}
  \lean{CombinatorialRigidity.Search.reachableFinding}
  \leanok
  \uses{def:directed-walk}
  Let $V$ be finite with decidable equality and $R$ a relation on $V$
  (with $R$ decidable in each argument, e.g.\ presented as an
  out-neighbour function $\mathrm{succ} : V \to \mathrm{List}\,V$,
  giving $R\,a\,b = (b \in \mathrm{succ}\,a)$). Given a source
  vertex $v$ and a predicate $P : V \to \mathrm{Bool}$, the
  procedure $\mathrm{reachableFinding}\,\mathrm{succ}\,P\,v$ performs
  a depth-first search along $R$-arcs from $v$, terminating at one of:
  \begin{itemize}
    \item (success) the first reachable vertex $w$ with $P(w)$,
      returning $\mathrm{some}\,(p, w)$ where $p$ is a simple
      $R$-walk from $v$ to $w$; or
    \item (failure) when the $R$-reachability closure of $v$
      contains no vertex satisfying $P$, returning $\mathrm{none}$.
  \end{itemize}
  Termination measure: $(\mathrm{Finset.univ} \setminus
  \mathrm{visited}).\mathrm{card}$, where $\mathrm{visited}$ is the
  finite set of vertices the DFS has already explored. Each
  recursive call expands $\mathrm{visited}$ by exactly one vertex,
  so the measure strictly decreases. (Compare
  \texttt{Batteries.UnionFind}'s
  \texttt{termination\_by self.rankMax - self.rank x}: the shape is
  the same --- a strictly-decreasing finite-domain measure --- but
  specialised to a $\mathrm{Finset}$-cardinality on a state set.)
\end{definition}

\begin{theorem}[Correctness of \texorpdfstring{$\mathrm{reachableFinding}$}{reachableFinding}]
  \label{thm:reachable-finding-correct}
  \lean{CombinatorialRigidity.Search.reachableFinding_sound,
    CombinatorialRigidity.Search.reachableFinding_complete}
  \leanok
  \uses{def:reachable-finding, def:directed-walk}
  Let $R$, $v$, $P$ be as in \cref{def:reachable-finding}. Then:
  \begin{enumerate}
    \item (soundness) If $\mathrm{reachableFinding}\,R\,v\,P =
      \mathrm{some}\,(p, w)$, then $p$ is a simple $R$-walk from $v$
      to $w$ and $P(w)$ holds.
    \item (completeness) If some vertex $w$ in the $R$-reachability
      closure of $v$ satisfies $P(w)$, then
      $\mathrm{reachableFinding}\,R\,v\,P = \mathrm{some}\,(p, w')$
      for some such $w'$ (possibly $w' \ne w$, since the DFS returns
      the first match in its traversal order, not necessarily the
      one the existence proof points at).
  \end{enumerate}
\end{theorem}
\begin{proof}
  \leanok
  \uses{def:reachable-finding}
  Both parts are by well-founded induction on the termination
  measure of \cref{def:reachable-finding}. Soundness: at each
  recursive step the walk extends by one arc (the DFS advances
  along an $R$-arc); the predicate is checked at each newly visited
  vertex before recursion; simplicity follows from the visited-set
  invariant (the DFS does not recurse into already-visited
  vertices). Completeness, by contradiction: assume some reachable
  $w$ satisfies $P$ and the DFS returns $\mathrm{none}$. The
  reachability witness lifts to an explicit $R$-walk $p$ from $v$
  to $w$ by head-first induction on $\mathrm{ReflTransGen}$;
  combining $p$ with the helper invariant
  \emph{``$\mathrm{aux}\,\mathrm{visited}\,v = \mathrm{none}$ implies
  no $R$-walk from $v$ avoiding $\mathrm{visited}$ ends at a
  $P$-satisfying vertex''} at $\mathrm{visited} = \varnothing$ then
  contradicts $P(w)$. The helper is proved by outer auto-induction
  on the DFS recursion (case-splitting on visited / $P\,v$ / else)
  with an inner induction on the walk-length bound: when the inner
  walk revisits the current source vertex, a sub-walk extracted via
  \texttt{DirectedWalk.dropUntilBundle} (truncation at first
  occurrence, bundled with its length-bound and vertex-subset
  witnesses) recurses by the inner IH; otherwise the walk-tail
  avoids the enlarged visited set and the outer IH closes the case.
\end{proof}

\emph{Reuse pattern.} The pebble game's \cref{def:tryReachPebble}
specialises this primitive: it instantiates $R$ to be the outgoing-arc
relation of the partial orientation, $P$ to be ``has a free pebble'',
and on success post-composes a path-reversal step that flips every
arc of the returned walk. The DFS itself is fully reusable across
this and any future application that needs reachability with a
witness path; the path-reversal is pebble-game-specific.

\subsection{Reachability closure (computable)}
\label{ssec:reach-closure-computable}

The pebble game's completeness side materialises a different view
of $\mathrm{succ}$-reachability from $v$: not the first match against
some predicate $P$, but the \emph{full set} of vertices reachable
from $v$, as a $\mathrm{Finset}\,V$. We ship this
view as a computable primitive
$\mathrm{reachClosureComputable}$, which avoids the
$\mathtt{Classical.decPred}$ dependency that a
filter-based formulation would introduce. Its semantic
contract --- characterisation as
$\mathrm{Relation.ReflTransGen}$ of the
$\mathrm{succ}$-arc relation --- is the consumer-facing iff
relied upon by the pebble-game's $D.\mathtt{reach}$ accessor in
\cref{sec:pebble-game}.

\begin{definition}[Reachability closure]
  \label{def:reachClosureComputable}
  \lean{CombinatorialRigidity.Search.reachClosureComputable}
  \leanok
  \uses{def:reachable-finding}
  Let $V$ be finite with decidable equality and
  $\mathrm{succ} : V \to \mathrm{List}\,V$ an out-neighbour function.
  Then $\mathrm{reachClosureComputable}\,\mathrm{succ}\,v :
  \mathrm{Finset}\,V$ is the set of vertices reachable from $v$
  along the relation $R\,a\,b = (b \in \mathrm{succ}\,a)$. The Lean
  implementation tests each candidate $w : V$ via
  $\mathrm{reachableFinding}\,\mathrm{succ}\,(\cdot = w)\,v$,
  collecting those for which the DFS succeeds; this routes through
  the iterative DFS already in scope rather than introducing a
  second iterative state machine.
\end{definition}

\begin{lemma}[Membership in the reach closure]
  \label{lem:mem-reachClosureComputable}
  \lean{CombinatorialRigidity.Search.mem_reachClosureComputable,
    CombinatorialRigidity.Search.reachClosureComputable_sound,
    CombinatorialRigidity.Search.reachClosureComputable_complete}
  \leanok
  \uses{def:reachClosureComputable, thm:reachable-finding-correct}
  For any $\mathrm{succ}$ and $v, w : V$,
  \[
    w \in \mathrm{reachClosureComputable}\,\mathrm{succ}\,v
      \;\Longleftrightarrow\;
      \mathrm{Relation.ReflTransGen}\,
        (\lambda a\,b \mapsto b \in \mathrm{succ}\,a)\,v\,w.
  \]
\end{lemma}
\begin{proof}
  \leanok
  \uses{def:reachClosureComputable, thm:reachable-finding-correct}
  Soundness: a successful DFS produces a witness walk by
  \cref{thm:reachable-finding-correct}; lifting that walk to
  $\mathrm{ReflTransGen}$ is a one-step induction
  ($\mathtt{DirectedWalk.toReflTransGen}$). Completeness:
  given the reachability chain, instantiate
  $\mathrm{reachableFinding}$ at the predicate $(\cdot = w)$ with
  $w$ as the existence witness; \cref{thm:reachable-finding-correct}
  returns $\mathtt{some}$, so $w$ passes the
  $\mathrm{Finset.filter}$ predicate.
\end{proof}
